\section{Related Work}
\label{sec:related}

\para{LLMs as recommendation models.} Early work showed that recommendation tasks can be cast as language generation. P5 reframed multiple recommendation problems as text-to-text prediction \cite{geng2022p5}, while instruction-following recommendation used templated instructions and task-specific tuning to improve recommendation quality \cite{zhang2023instruction}. Survey papers argue that LLMs add semantic reasoning, flexible interfaces, and stronger zero-shot behavior, but they also emphasize unresolved issues around grounding, evaluation, and personalization \cite{wu2023survey,zhao2023era}. Our work is narrower: we do not tune a model or train an adapter, and we focus on whether prompting alone is competitive on public music ranking tasks.

\para{Prompting versus hybrid recommendation.} Several recent systems suggest that LLMs are often strongest when paired with classical recommender components rather than used alone. LLM-Rec uses prompting to enrich item text for downstream recommenders \cite{lyu2023llmrec}. \textsc{LLaRA} aligns sequential recommender embeddings with LLM inputs \cite{liao2024llara}, and \textsc{A-LLMRec} injects collaborative knowledge from a frozen collaborative filter into a frozen LLM \cite{kim2024allmrec}. These systems report strong gains, but they answer a different question from ours because they rely on learned recommenders, alignment modules, or training data beyond the prompt itself. We instead isolate the prompt-only regime.

\para{Music recommendation and playlist continuation.} Public music recommendation benchmarks remain centered on playlist continuation and next-item prediction. The \mpd challenge literature provides a strong non-LLM reference point for Spotify-style recommendation, including collaborative filtering, neighbor methods, and hybrid reranking systems \cite{teinemaa2018playlist,zamani2018analysis}. More recently, \textsc{TalkPlay} shows that LLM-based music recommendation can work well when the model is trained end to end over multimodal music representations \cite{doh2025talkplay}. Our work sits between these lines of research: unlike the playlist continuation literature, we evaluate frontier LLMs directly; unlike \textsc{TalkPlay}, we do not train a music-specific model.

\para{Popularity bias and beyond-accuracy evaluation.} Recommendation quality is not only about top-1 or top-10 accuracy. Recent work studies how LLM recommenders interact with popularity bias and catalog concentration \cite{lichtenberg2024popularity}. That perspective is especially relevant for music, where users often value novelty and stylistic fit. Following this line of work, we report both ranking metrics and simple beyond-accuracy statistics, including mean top-10 popularity percentile and coverage.

Taken together, prior work shows that LLMs can contribute meaningfully to recommendation, especially in hybrid systems. The missing piece is a clean measurement of prompt-only recommendation quality in a music setting using the same candidate sets as classical baselines. That is the gap we target.
